\scanpage{421}
\mspage{213}
% Main text, approximate manuscript lines 1--4.
\deleted{Soit} L'unicité est-elle de \unsure{ttes|toutes} façons claire ?
\deleted{\unsure{Choisissant un}} Considérons l'ext. \inserted{$\mathcal{E}$} de
$\Gamma$ par $L=\mathrm{Ker}((N_a)_{I'}\to\unsure{N_a|N_0})$,
\unsure{image} inverse de l'ext. $(N_a)_{I'}$
(de $\unsure{N_{a'}|N_{0'}}$ par $L$) via $\Gamma\to\unsure{N_{a'}|N_{0'}}$, on
% Main text, approximate manuscript lines 5--8.
a un scindage \inserted{partiel} de cette extension \unsure{au} dessus
du sous-groupe $\Gamma'$ de $\Gamma$, et l'assertion est
que ce scindage se \uline{\unsure{prolonge}}, \uline{de façon} \uline{unique}, à $\Gamma$.
OPS $\Gamma'$ invariant dans $\Gamma$, et on
% Main text, approximate manuscript lines 9--13.
identifie $\Gamma'$ à un ss-groupe de $\mathcal{E}$. \unsure{Toutes}
\deleted{\unsure{sections}} les sections de $\mathcal{E}$ sur $\Gamma$ \unsure{s'identifient}
à des sous-groupes \inserted{\unsure{sections}} $\widetilde{\Gamma}$ de $\mathcal{E}$.
Pour un tel $\widetilde{\Gamma}$, on a bien sûr
$\widetilde{\Gamma}\subset\mathrm{Norm}_{\mathcal{E}}(\Gamma')$
d'ailleurs \deleted{\unsure{si on} \illegible{2}} \inserted{on a évidemment}
% Main text, approximate manuscript line 14.
\[
\mathrm{Norm}_{\mathcal{E}}(\Gamma')\cap L=L^{\Gamma'} \tag{31}
\]
% Main text, approximate manuscript lines 15--21.
et je présume qu'on doit avoir
$L^{\Gamma'}=1$ (qui généralise la condition
\textcircled{3} plus haut\ldots) Or \uline{cette condition}
\uline{implique l'unicité} -- savoir
$\widetilde{\Gamma}=\mathrm{Norm}_{\mathcal{E}}(\Gamma')$ --
et l'existence \uline{signifie} que
$\mathrm{Norm}_{\mathcal{E}}(\Gamma')\to\Gamma$
est épi (donc iso\ldots).

% Main text, approximate manuscript lines 22--23.
Mais il faudra essayer de préciser,
dans certains contextes (au moins celui des
\scanpage{422}
% Main text, approximate manuscript lines 1--4; continuation of scan 421.
actions arithmétiquement fidèles, i.e. $\Gamma\to\boldsymbol{\Gamma}_a$
injectif -- qui \unsure{correspond} \unsure{normalement|vraisblmt}
au cas de courbes alg. \unsure{définies} sur des
ext. algébriques de $\mathbb{Q}$ -- la notion
% Main text, approximate manuscript lines 5--9.
d'action ``admissible''. Pour ceci, on
doit revenir sur la relation entre
courbes (potentielles) et revêtements finis,
ce qui donne aussi une façon de
faire varier $g$.

% Main text, approximate manuscript lines 10--17.
Mais j'ai envie d'abord de reprendre
\unsure{sous|sur} un autre aspect \inserted{(peut-être plus général)} le formalisme
précédent, qui \unsure{peut} peut-être aussi \unsure{s'}appliquer au cas des groupes \uline{discrets}
à lacets (je pense au formalisme
des $\unsure{U^!|\mathcal{U}^!}$, lié aux actions extérieures de
groupes finis sur de tels \unsure{gr.-p.} à lacets).

% Main text, approximate manuscript lines 18--20.
Soit $\mathcal{X}$ un ensemble \inserted{$\ne\varnothing$}
(\inserted{\unsure{moralement}, un} ensemble des
``pts'' d'une \deleted{surface} courbe algébrique,
ou d'une surface topologique\ldots), \deleted{et $\Gamma$}
\scanpage{423}
\mspage{214}
% Main text, approximate manuscript lines 1--7.
\deleted{un groupe opérant sur $\mathcal{X}$.} On suppose
donné, pour tte partie finie $I$ de $\mathcal{X}$,
de complémentaire $U_I=\mathcal{X}\setminus I$, un
\unsure{groupoïde} $\Pi_{U_I}$, ayant $U_I$ comme ens
d'\unsure{objets}. De plus, pour $J\supset I$ i.e.
$U_J\subset U_I$, on suppose donné un hom
de groupoïdes
% Main text, approximate manuscript line 8.
\[
\Pi_{U_J}\to\Pi_{U_I}
\]
% Main text, approximate manuscript lines 9--14.
qui sur les objets soit l'inclusion $U_J\unsure{\subset|\hookrightarrow}\unsure{U_J|U_I}$.
On suppose transitivité (stricte) i.e. \deleted{\unsure{qui}}
un foncteur covariant \inserted{de la cat.} des parties $U$ de
$\mathcal{X}$ complémentaires de parties finies \inserted{(avec les inclusions)}, vers
\unsure{la} cat. des groupoïdes, qui sur \unsure{l'ens.}
d'objets coïncide avec le foncteur évident\ldots

% Left marginal note, running beside approximate lines 10--22.
\margin{Si $\mathcal{X}$ est fini et $I=\mathcal{X}$, on
\unsure{suppose} \inserted{\uuline{\unsure{quand même|cependant}}}
$\Pi_{U_I}\;(=\Pi_{\varnothing})$ \unsure{n'est}
pas le groupoïde vide, mais un groupoïde
\deleted{quelconque} (qui \unsure{s'envoie} dans les précédents\ldots)
mais on ne pourra exiger
\unsure{la} transitivité stricte.
Il \unsure{faudra|faudrait} \inserted{\unsure{peut-être}} \illegible{1} \unsure{hyp.} anabéliennes
sur $\Pi_{U_I}$ dans ce cas\ldots}

% Main text, approximate manuscript lines 15--18.
On suppose de plus que pour tout \deleted{$I\in\mathcal{P}_f(\mathcal{X})$}
$i\in\mathcal{X}$, on se donne un \deleted{\unsure{groupoïde} \illegible{1}}
$\Pi_i$ (ou $\Pi_{D_i^*}$),
et \deleted{\unsure{pour tout}} pour $i\in I\in\mathcal{P}_f(\mathcal{X})$,
un hom de groupoïdes
% Main text, approximate manuscript line 19.
\[
\Pi_{D_i^*}\to\Pi_{U_I}
\]
% Main text, approximate manuscript lines 20--23.
compatible avec les morphismes de transition
\inserted{$\Pi_{U_I}\to\Pi_{U_J}$ $(I\supset J)$}.
On suppose que pour $I$ fixé, on trouve
sur $\Pi_{U_I}$ une structure de groupoïde
à lacets par
\[
\Pi_{D_I^*}=\coprod_{i\in I}\Pi_{D_i^*}\to\Pi_{U_I}.
\]
\scanpage{424}
% Main-text lines 1--2; insertion above line 2.
sans d'ailleurs exclure le cas où $I=\varnothing$, et où \inserted{($\Pi_{\mathcal{X}}$ \unsure{étant} connexe, disons)} \unsure{son|le} genre \uline{est $0$} --- \deleted{\unsure{donc}}
\margin{\unsure{On suppose} \illegible{4} \unsure{que} \illegible{1} $\mathcal{X}$ \illegible{3} \deleted{\unsure{pour} $\mathcal{X}$ \illegible{2}} \deleted{\unsure{deux pts}}}

% Main-text lines 3--8.
On suppose de plus que si $J\supset I$ i.e. $U_J\subset U_I$, \inserted{et pour $i\in J\setminus I$,} l'hom. composé
\[
\Pi_{D_i^*}\longrightarrow\Pi_{U_J}\longrightarrow\Pi_{U_I}
\]
soit ``constant'', de telle façon que $\Pi_{U_I}$ se \unsure{déduise} de $\Pi_{U_J}$ par ``bouchage des trous'' en $J\setminus I$.

% Main-text lines 9--21; the middle of the example is heavily revised.
Jusqu'à présent, \unsure{tt|tout} ceci pouvait se visualiser p.ex. en partant d'une surface top. orientable \deleted{de la forme \unsure{$X\setminus S$}; \illegible{4} du type} $X$, et d'une partie \inserted{\illegible{1}} $\mathcal{X}$ de $X$ \deleted{\illegible{1}}\inserted{\unsure{rencontrant}} \deleted{\unsure{tte compos. connexe}}\inserted{\unsure{tte} comp. connexe}; en prenant pour $\Pi_I$ la restriction à $\mathcal{X}\setminus I$ du groupoïde fondamental de $X\setminus I$ (si $I\ne\mathcal{X}$; si $\mathcal{X}=I$, ce qui exige $\mathcal{X}$ fini, on prendra le groupoïde fondamental de $X\setminus I=X\setminus\{\mathcal{X}\}$; qu'on aura du mal à envoyer dans les autres avec transitivité \uline{stricte}, qu'à cela ne tienne !

% Main-text lines 22--23; sentence continues on the next scan.
Mais, on est surtout intéressé au cas $\mathcal{X}$ infini. Où on \unsure{prend} $X$ courbe alg. sur
\scanpage{425}
\mspage{215}
% Main-text lines 1--3; continuation from scan 424.
corps alg. clos, $\mathcal{X}$ partie de $X$ rencontrant \unsure{tte} comp. connexe, et on définit les $\Pi_{U_I}$ comme précédemment.

% Main-text lines 4--9, with an insertion between lines 6 and 7.
Soit maintenant $\Gamma$ un groupe qui opère sur la situation (\uline{strictement}, s'entend) donc sur \unsure{l'ens} $\mathcal{X}$, \deleted{et} les groupoïdes \unsure{$\Pi_{U_I}$}, \inserted{les $\Pi_{D_i^*}$, (N.B. qu'il \unsure{permute} entre eux),} en commutant aux hom. \unsure{d|de}
\[
\Pi_{U_I}\longrightarrow\Pi_{U_J}
\]
et les $\Pi_{D_i^*}\longrightarrow\Pi_{U_I}$.
\margin{\deleted{(qu'il \unsure{permute} \unsure{entre eux})}}

% Main-text lines 10--13; the wording after Nous is rewritten interlinearly.
(On aura des difficultés, si $\mathcal{X}$ est fini, pour $\Pi_{U_I}$ quand $I=\mathcal{X}$, pour la commutation \uline{stricte} \deleted{\unsure{des}} \unsure{$\Pi_{U_{\mathcal{X}}}$}$\longrightarrow\Pi_{U_J}$ \unsure{..} mais passons...)
Nous \deleted{faisons \unsure{l'hypothèse|les hypothèses},} \inserted{supposons définie une notion de ss-groupe \uline{admissible} de $\Gamma$, (ils sont en tous cas $\ne1$) telle que \deleted{\illegible{1}} si $\Gamma'$ est admissible \unsure{alors} tt ss-groupe $\Gamma''\supset\Gamma'$ aussi.} \deleted{\unsure{suivantes}} \inserted{On \deleted{\unsure{suppose}}}

% Canceled condition and canceled bracketed development in the lower half.
\deleted{\textcircled{\(1^\circ\)} Pour $P\in\mathcal{X}$, son stabilisateur $\Gamma_P$ dans $\Gamma$ est \deleted{$\ne1$}\inserted{admissible} (donc $\ne1$)}
\inserted{\deleted{\unsure{$U_I$ stable pr $\Gamma_P$ i.e. $I$ stable}}}

\deleted{\unsure{\textcircled{\(2^\circ\)}} \illegible{1} (Notons que si $U_I\ni P$, alors $P$ est un objet de $\Pi_{U_I}$ avec op. de $\Gamma_P$ dessus, \deleted{\illegible{1}}\inserted{On} \unsure{trouve} ainsi une \deleted{application}\inserted{opération} de $\Gamma_P$ sur \unsure{$\pi_1(U_I,P)$|$\Pi_1(U_I,P)$}, et on suppose}

\deleted{\textcircled{\(2^\circ\)} \unsure{$\pi_1(U_I,P)$|$\Pi_1(U_I,P)$} \deleted{\unsure{appeler}}\inserted{\unsure{dire qu'un}} ss-groupe $\Gamma'$ de $\Gamma$ est \textcircled{\(1^\circ\)} Si $\Gamma'$ \uline{admissible} si \deleted{\illegible{1}} $\mathcal{X}^{\Gamma'}$ \unsure{rencontre} \unsure{tte} composante connexe de $\mathcal{X}=U_{\varnothing}$ (ce qui}
\scanpage{426}
% Main-text lines 1--3 continue the canceled draft on scan 425.
\deleted{signifie aussi \unsure{simplement} que $\mathcal{X}^{\Gamma'}\ne\varnothing$ si $\mathcal{X}$ connexe, ou encore que $\Gamma'$ \unsure{contient} un stabilisateur $\Gamma_P$.} On suppose

% Main-text line 4 and its small parenthetical addition.
\textcircled{\(1^\circ\)} $\forall P\in\mathcal{X}$, $\Gamma_P\ne\{1\}$, et $\Gamma_P$ d'indice fini \unsure{ds $\Gamma$} \inserted{(i.e. l'orbite de $P$ finie)}

% Main-text lines 5--21.
Nous voulons des conditions qui assurent que $\mathcal{X}$ et les $\Pi_{U_I}$ \deleted{\unsure{pour}} etc, \unsure{peuvent|puissent} se reconstruire à \deleted{\unsure{l'aide}} \unsure{isom.} can. près à l'aide des données purement groupoïdiques ``nues'' ou topossiques \inserted{(à opérateurs $\Gamma$)} \unsure{correspondants|correspondantes}. \deleted{On \illegible{1}} On peut le dire en langage topossique savant, mais on va l'exprimer en termes de théorie des groupes à lacets. Soit $I$ une partie de $\mathcal{X}$ stable par \unsure{$\Gamma^*$|$\Gamma'$}, on voudrait poser des conditions qui \deleted{\unsure{on prend} un \unsure{ss-groupe}} \unsure{permettent|permettront} d'exprimer \inserted{\unsure{(pour les tels $I$)}} la situation des groupoïdes \inserted{$\Pi_V$} associés aux $V\supset U_I$ (i.e. les $U_J$ avec $J\subset I$), \inserted{les $\Pi_{D_i^*}$ $(i\in I)$,} et l'opération de $\Gamma$ dessus, en termes de la \deleted{\unsure{op.}} seule opération de $\Gamma$ sur le groupe ext. à lacets \deleted{$\pi_I$ associé} $\pi_1(\Pi_{U_I})$ associé à $I$ ---
\scanpage{427}
\mspage{216}
\margin{\unsure{$\rightsquigarrow$|$\longrightarrow$}}
% Main-text lines 1--5.
ou plutôt (car cela est immédiat, par l'opération de forage des trous), \inserted{en sens inverse,} montrer comment,
\deleted{\unsure{pour un} $J\subset I$ \unsure{invariant} par $\Gamma$}
\deleted{\illegible{5}}
\deleted{\unsure{réconstituer} \illegible{2},}
% Main-text lines 6--11.
sous réserve d'anabélianité \unsure{$\zeta$|$\xi$} disons de \unsure{$\Pi_{\mathcal X}$|$B_{\Pi_{\mathcal X}}$} (pour fixer les idées), on peut $\pm$ \unsure{réconstituer} $\Pi_{U_I}$ et l'action de $\Gamma$ dessus \unsure{....|...},
à partir de \unsure{$\Pi_{\mathcal X}$} \inserted{(ou plutôt, \deleted{de}\inserted{du \unsure{top.}} \unsure{$B_{\Pi_{\mathcal X}}$})} et de l'action de $\Gamma$ dessus.
On veut au moins une description \inserted{\unsure{intrinsèque}} des \unsure{éléments} de \unsure{$\mathcal X$}, \inserted{et de l'action de $\Gamma$ dessus,} via
% Main-text lines 12--15.
cette action ``molle'' --- qui, pour $X$ connexe, revient encore à une action extérieure de $\Gamma$ sur un groupe à lacets (\uline{sans} lacets!) $\pi$...

% Main-text lines 16--21; the parenthesis continues on the next scan.
Une première condition, sous forme faible, \unsure{c'est|est} que (supposant \unsure{$\mathcal X$} connexe, et appelant : $\pi(I)$ le groupe extérieur $\pi_1(\Pi_{U_I})$, pour \unsure{tte} partie finie $I$ de \unsure{$\mathcal X$}) que \unsure{pour} $\pi(I)$ anabélien (ce qui sera le cas pour $I$ ``assez grand'',
\scanpage{428}
% Main-text lines 1--3.
du moins si \unsure{$\mathcal X$} est infini \unsure{donc|de} qu'on puisse prendre $I$ de cardinal arbitrairement grand...) \deleted{on ait} l'application évidente
% Deleted display and replacement display, approximately lines 4--5.
\[
\deleted{\unsure{\mathrm{Pt}(B_{\pi(I),\Gamma\natural})|\mathrm{Pt}(B_{\pi(I)},\Gamma\natural)}\longrightarrow\text{classes de}}
\]
\[
2^\circ)\quad
\deleted{\unsure{\mathcal X}}U_I\longrightarrow
\unsure{\mathrm{Pt}(B_{\pi(I),\Gamma\natural})|\mathrm{Pt}(B_{\pi(I)},\Gamma\natural)}
\simeq
\begin{gathered}
\unsure{\inserted{\text{classes de }\pi(I)\text{-conj. }\unsure{\text{de}|\text{des}}}}\\
\unsure{\text{germes}}\text{ de}\\
\deleted{\text{de}}\text{ scindages de}\\
\text{l'ext. }\mathcal E(I)\\
\text{de }\deleted{\unsure{\mathcal E}}\Gamma\text{ par }\pi(I)
\end{gathered}
\]
% Main-text lines 6--15.
\unsure{(compatible|compatible} avec les actions de $\Gamma$) soit \uline{injective}. Mais pour aller plus loin, il faudrait pouvoir donner une caractérisation de l'image. Notons (nous plaçant dans le contexte profini désormais) que si on se donne sur les \unsure{$\Pi(I)$|$\pi(I)$} des \unsure{arithmétisations}, \unsure{compatibles} avec les hom extérieurs de bouchage de trous $\pi(I)\longrightarrow\pi(J)$, et avec l'action de $\Gamma$ sur les $\pi(I)$, on a donc un
% Main-text lines 16--20.
hom canonique de $\Gamma$ dans le groupe commun \deleted{des} $\boldsymbol{\Gamma}$ des autom. \unsure{arithmétiquement} \unsure{(}ext. de ces \unsure{arithmétisations}.
(et $\overline{\mathrm m}$ $\Gamma\longrightarrow$ \unsure{$M_{\pi(\varnothing)=\pi}$}, --- qui est \inserted{déjà} une \inserted{donnée} nettement plus forte). On peut donc
\scanpage{429}
\mspage{217}
% Main text, approximate lines 1--5.
supposer \unsure{donné} une notion de \uline{relèvement}
\uuline{admissible} d'une telle action arithm.
extérieure, de telle façon que le
\unsure{mor.} précédent soit une bijection
\[
 U_I=\unsure{\mathcal X|\mathfrak X}\setminus I
 \xrightarrow{\sim}
 \mathrm{Pt}_{\mathrm{ad}}\bigl(B_{\pi(I)},\Gamma\natural\bigr)
\]
Ceci signifie d'ailleurs, pratiquement,
que la notion d'admissibilité satisfait
aux conditions $1^\circ)$ : $6^\circ)$ \unsure{vues} ci-dessus.

% Main text, approximate lines 9--12.
Quant à savoir \unsure{ce} qu'il y a lieu
d'appeler \unsure{op} \deleted{\unsure{ext.}} ``\unsure{admissibles|admissible}'' de
$\Gamma$ sur un groupe \inserted{ext.} \unsure{(}profini\unsure{)|,} à lacets,
cela reste pour le moment conjectural.

% Main text, approximate lines 13--21; interlinear revisions retained.
On pourrait à titre expérimental
\unsure{conjecturer} \deleted{que \unsure{toute} \illegible{2}}
\inserted{que la notion qui suit marcherait}
\unsure{Appelons|Appellons}
``admissible'' toute telle action
\[
 \Gamma\longrightarrow\widehat{\widehat{\mathcal T}}(\pi)
\]
qui respecte une arithmétisation $a$,
i.e. qui se factorise par le $N_a$ correspondant, \unsure{et} telle que l'im. de
$\Gamma\longrightarrow N_a/\widehat{\mathcal T}_a=\boldsymbol{\Gamma}_a$
soit \unsure{ouverte|ouvert},
\deleted{\inserted{et que $\pi^{\Gamma'}=(1)$}}
\deleted{\unsure{les conditions}}, \inserted{et que}
\scanpage{430}
% Main text, approximate lines 1--5.
\deleted{\unsure{envisagées} plus haut \unsure{$1^\circ:6^\circ$} \unsure{sont} \unsure{vérifiées|vérifiés}}

$\forall$ \unsure{hom} de bouchage de trous en $i\in I(\pi)$,
$\pi\longrightarrow\pi_i'$, $N(\pi)_i\longrightarrow M(\pi_i')$, l'hom \unsure{correspondant}
$\Gamma_i$, ni d'aucun sous-groupe ouvert $\Gamma'$
de $\Gamma_i$, \unsure{ne} normalise un $L_j'$ $(j\in I(\pi_i')=I\setminus\{i\})$.

% Upper left margin; several long strokes run through the whole note.
\margin{\deleted{
cette condition est
peut-être \illegible{2} et \unsure{prématurée}
--- il se pourrait
que pour \unsure{tout couple|tous couples}
d'indices distincts $i,j$ de $I=I(\pi)$,
\inserted{et pour $L_i,L_j$ \unsure{correspondants}},
le sous-groupe \deleted{\illegible{1}} \unsure{$K$} de
$M(\pi)$ \unsure{formé} des
$g\in M(\pi)$ qui \uline{normalisent}
$L_i$, et qui normalisent de
plus $L_j$ \uline{mod} $\mathrm{Ker}(\unsure{\pi\longrightarrow\pi_i'})$,
soit ``\unsure{trop petit}'' \deleted{\unsure{pour que}} \inserted{\unsure{i.e.}}
que son image dans
$\unsure{\boldsymbol{\Gamma}|\Pi}=M(\pi)/\unsure{\widehat{\mathcal T}(\pi)|\widehat{\mathcal G}(\pi)}$
\uline{\unsure{ne soit pas}}
d'indice fini
}}

% Main text, approximate lines 6--14.
\unsure{peut-être} $\overline{\mathrm m}$ faudrait-il imposer
$\pi^{\unsure{(\Gamma^{\prime+})|\langle\Gamma^{\prime+}\rangle}}=(1)$, où $\Gamma^{\prime+}$ est le noyau de
$\Gamma'\xrightarrow{\chi}\widehat{\mathbb Z}^{*}$ --- en tous cas on s'attend
\unsure{---} que l'on ait alors $\pi^{\Gamma'}=(1)$,
et si ce n'était le cas, il faudrait
l'introduire dans la définition --- pour
chaque opération de bouchage de trous
\unsure{relatif} à $I'\subset I(\pi)$, et au choix d'un
$i\in I'$, permettant de définir $\Gamma_{(I',i)}\longrightarrow M(\pi(I'))$.

% Lower left marginal note.
\margin{
Plus généralement,
pour $i\in I'\subset I=I(\pi)$,
\deleted{\unsure{il faut}} \inserted{\unsure{on impose}} que la
vraie action de
\unsure{$\Gamma_{I',i}$} sur \unsure{$\pi_i(I')$|$\Pi_i(I')$}
\inserted{\unsure{correspondante}} \unsure{à $i$}
ne normalise aucun
sous-groupe $L_j$
$(j\in I\setminus I')$ --- on
est ramené pour
ceci au cas où
$I'=I\setminus\{j\}$, donc
où \unsure{$\pi(I')$|$\Pi(I')$} n'a \unsure{(qu'une} seule \unsure{classe} lacets\dots\
(\unsure{si|et} $g\geq1$)
}

% Main text, approximate lines 15--21.
Notons qu'une vraie action de $\Gamma$ sur $\pi$
(pas \unsure{seult|seulmt} \unsure{ext.}) qui relève une action
donnée\unsure{$+$} admissible, est \inserted{\unsure{elle aussi}} admissible
(en ce sens qu'elle définit une
action \deleted{ext.} \uline{adm} sur un $\pi'$ de
type $(g,\nu+1)$), ssi cette action \inserted{\unsure{(plutôt son germe)}} ne
normalise aucun $L_i$ $(i\in I(\pi))$, et qu'
